\section{Matriz de cambio de base}
\begin{defin}
Sean $B_1 = \{u_1 , u_2, \dots u_n\}$ y $B_2 = \{ v_1,v_2, \dots, v_n \}$ dos bases en $\mathbb{R}^n$. Sea la matriz 
$$A = \begin{pmatrix} a_{11} & a_{12} & a_{13} & \cdots & a_{1n} \\ a_{21} & a_{22} & a_{23} & \cdots & a_{2n} \\ \vdots & \vdots & \vdots & & \vdots \\ a_{n1} & a_{n2} & a_{n3} & \cdots & a_{nn} \end{pmatrix}$$

 donde:
\begin{center}
     $(u_1)_{B_2} = \begin{pmatrix} a_{11} \\ a_{21} \\ \vdots \\ a_{n1} \end{pmatrix}$, $(u_2)_{B_2} = \begin{pmatrix} a_{1n} \\ a_{2n} \\ \vdots \\ a_{nn} \end{pmatrix}$ $\cdots, (u_n)_{B_2} = \begin{pmatrix} a_{1n} \\ a_{2n} \\ \vdots \\ a_{nn} \end{pmatrix}$. 
\end{center}
La matriz $A$ se denomina la matriz de transición de la base $B_1$ a la base $B_2$.
\end{defin}




\begin{thm}
Sean $B_1$ y $B_2$ bases para un espacio vectorial $V$. Sea $A$ la matriz de transición de $B_1$ a $B_2$. Entonces para todo $x \in V$ $(x)_{B_2} = A(x)_{B_1} $.
\end{thm}

\begin{thm}
Sea $A$ la matriz de transición de $B_1$ a $B_2$. Entonces $A^{-1}$ es la matriz de transición de $B_2$ a $B_1$.
\end{thm}

\begin{itemize}
\item
\textbf{Procedimiento} para encontrar la matriz de transición de la base canónica a la base $B_2 = \{ v_1,v_2, \dots v_n \}$:

\begin{enumerate}
\item
Se escribe la matriz $C$ cuyas columnas son $v_1,v_2, \dots, v_n$.

\item
Se calcula $C^{-1}$. Esta es la matriz de transición que se busca.
\end{enumerate}
\end{itemize}





\begin{ejem}
\begin{enumerate}
\item
En $\mathbb{R}^3$, sea $B_1 = \left\{ \begin{pmatrix} 1\\ 0 \\ 0  \end{pmatrix} ,\begin{pmatrix} 0 \\ 1 \\ 0  \end{pmatrix}, \begin{pmatrix} 0\\ 0 \\ 1  \end{pmatrix} \right\}$ y $B_2 = \left\{ \begin{pmatrix} 1\\ 0 \\ 2  \end{pmatrix} ,\begin{pmatrix} 3 \\ -1 \\ 0  \end{pmatrix}, \begin{pmatrix} 0\\ 1 \\ -2  \end{pmatrix} \right\}$. Si $x = \begin{pmatrix} x \\ y \\ z \end{pmatrix} \in \mathbb{R}^3$, encuentre $(x)_{B_2}$
\\\

\item
En $P_2$ la base canónics es $B_1 = \{ 1,x,x^2 \}$. Sea $B_2 = \{ 4x-1,2x^2-x,3x^2+3 \}$. Si $p = a_0 + a_1x + a_2x^2$. Encuentre $B_2$
\\\

\item
Sean $B_1 = \left\{ \begin{pmatrix} 3 \\ 1 \end{pmatrix} ,\begin{pmatrix} 2 \\ -1  \end{pmatrix} \right\}$ y $B_2 = \left\{ \begin{pmatrix} 2 \\ 4 \end{pmatrix} ,\begin{pmatrix} -5 \\ 3 \end{pmatrix} \right\}$. Encuentre la matriz de transición entre $B_1$ y $B_2$.
\end{enumerate}
\end{ejem}





\begin{ejem}
Sean $B_1 = \left\{ \begin{pmatrix} 3 \\ 1 \end{pmatrix} ,\begin{pmatrix} 2 \\ -1  \end{pmatrix} \right\}$, $B_2 = \left\{ \begin{pmatrix} 2 \\ 4 \end{pmatrix} ,\begin{pmatrix} -5 \\ 3 \end{pmatrix} \right\}$ y $C$ la base canónica. 

\begin{enumerate}
\item
Encuentre la matriz de transición de $C$ a $B_2$. Llame a esa matriz $A_1$

\item
Encuentre la matriz de transición de $B_1$ a $C$. Llame a esa matriz $A_2$

\item
Calcule $A_2A_1$. ¿Es este resultado el mismo que el resultado del último ejemplo?
\end{enumerate}
\end{ejem}





\begin{thm}
Sea $B_1 = \left\{ v_1,v_2, \dots, v_n \right\}$ una base del espacio vectorial $V$ de dimensión $n$. Suponga que 
\\\

$(x_1)_{B_1} = \begin{pmatrix} a_{11} \\ a_{21} \\ \vdots \\ a_{n1} \end{pmatrix}, (x_2)_{B_2} = \begin{pmatrix} a_{12} \\ a_{22} \\ \vdots \\ a_{n2} \end{pmatrix}, \cdots , (x_n)_{B_1} = \begin{pmatrix} a_{1n} \\ a_{2n} \\ \vdots \\ a_{nn} \end{pmatrix}$.
\\\

Sea $A = \begin{pmatrix} a_{11} & a_{12} & a_{13} & \cdots & a_{1n} \\ a_{21} & a_{22} & a_{23} & \cdots & a_{2n} \\ \vdots & \vdots & \vdots & & \vdots \\ a_{n1} & a_{n2} & a_{n3} & \cdots & a_{nn} \end{pmatrix}$.
\\\

Entonces $x_1, x_2, \dots, x_n$ son linealmente independientes si y solo si $\det A \not= 0$.
\end{thm}





\begin{ejem}
\begin{enumerate}
\item
En $\mathbb{P}_2$, determine si los polinomios $3-x; 2+x^2$ y $4+5x-2x^2$ son linealmente dependientes o independientes.
\\\

\item
En $M_{22}$ sea $B_1 = \left\{ \begin{pmatrix}  1 & 0 \\ 0 & 0 \end{pmatrix}, \begin{pmatrix}  0 & 1 \\ 0 & 0 \end{pmatrix}, \begin{pmatrix}  0 & 0 \\ 1 & 0 \end{pmatrix}, \begin{pmatrix}  0 & 0 \\ 0 & 1 \end{pmatrix} \right\}$ y 

$B_2 = \left\{ \begin{pmatrix}  1 & 0 \\ 0 & 0 \end{pmatrix}, \begin{pmatrix}  1 & 1 \\ 0 & 0 \end{pmatrix}, \begin{pmatrix}  1 & 1 \\ 1 & 0 \end{pmatrix}, \begin{pmatrix}  1 & 1 \\ 1 & 1 \end{pmatrix} \right\}$.Encuentre la matriz de transición de $B_1$ a $B_2$
\end{enumerate}
\end{ejem}





\subsection{Rango, Nulidad, Espacio Fila y Espacio Columna}

\begin{defin}
$N_A = \{ x \in \mathbb{R}^n:Ax =0 \}$ se denomina el espacio nulo de $A$ y $\nu (A)$=dim$N_A$ se denomina nulidad de $A$. Si $N_A$ contiene sólo al vector cero, entonces $\nu (A) = 0$.
\end{defin}

\begin{ejem}
\begin{enumerate}
    \item 
    Sea $A = \begin{pmatrix} 1 & 2 & -1 \\ 2 & -1 & 3 \end{pmatrix}$. Encuentre $N_A$.
\\\

\item
Sea $A = \begin{pmatrix} 2 & -1 & 3 \\ 4 & -2 & 6 \\ -6 & 3 & -9 \end{pmatrix}$. Encuentre $N_A$.
\end{enumerate}
\end{ejem}

\begin{thm}
Sea $A$ una matriz de $n \times n$. Entonces $A$ es invertible si y sólo si $\nu (A) = 0$.
\end{thm}



\begin{defin}
Sea $A$ una matriz de $m \times n$. Entonces la imagen de $A$, denotada por im$A$ está dada por im$A = \{ y \in \mathbb{R}^m : Ax = y$ para alguna $x \in \mathbb{R}^n \}$.
\end{defin}

\begin{thm}
Sea $A$ una matriz de $m \times n$. Entonces im$A$ es un subespacio de $\mathbb{R}^m$.
\end{thm}

\begin{defin}
Sea $A$ una matriz de $m \times n$. Entonces el rango de $A$, denotado por $\rho (A)$, está dado por $\rho (A) = $dim im$A$.
\end{defin}





\begin{defin}
Si $A$ es una matriz de $m \times n$, sean $\{ r_1, r_2, \dots, r_m \} $ las filas de $A$ y $\{ c_1, c_2, \dots, c_n \}$ las columnas de $A$. Entonces se define: 
\begin{itemize}
    \item 
    $R_A = $Espacio fila de $A =$ gen$ \{ r_1, r_2, \dots, r_m \}$.
    
    \item
     $C_A =$ espacio columna de $A = $ gen$\{ c_1,c_2, \dots c_n \}$. 
\end{itemize}
\end{defin}
    
    \begin{thm}
    Para cualquier matriz $A$, $C_A =$ im$A$.
    \end{thm}
    
    \begin{ejem}
    Sea $A = \begin{pmatrix} 2 & -1 & 3 \\ 4 & -2 & 6 \\ -6 & 3 & -9 \end{pmatrix}$. Encuentre $N_A$, $\nu(A)$, im$A$, $\rho(A)$, $R_A$, $C_A$.
    \end{ejem}
    






\begin{thm}
Si $A$ es una matriz de $m \times n$, entonces dim$R_A=$ dim$C_A=$ dim im$A = \rho (A)$.
\end{thm}

\begin{ejem}
Sea $A = \begin{pmatrix} 2 & -1 & 3 \\ 4 & -2 & 6 \\ -6 & 3 & -9 \end{pmatrix}$. Encuentre im$A$ y $\rho(A)$.
\end{ejem}

\begin{thm}
Si $A$ es equivalente por filas a $B$, entonces $R_A = R_B$, $\rho (A) = \rho (B)$ y $\nu (A) = \nu (B)$.
\end{thm}





\begin{thm}
El rango de una matriz es igual al número de pivotes en su forma escalonada por filas.
\end{thm}

\begin{ejem}
\begin{enumerate}
\item
Determine el rango y espacio fila de $A = \begin{pmatrix} 1 & -1 & 3 \\ 2 & 0 & 4 \\ -1 & -3 & 1 \end{pmatrix}$. 
\\\

\item
Encuentre una base para el espacio generado por $v_1 = \begin{pmatrix} 1 \\ 2 \\ -3 \end{pmatrix}, v_2 = \begin{pmatrix} -2 \\ 0 \\ 4 \end{pmatrix}, v_3 = \begin{pmatrix} 0 \\ 4 \\ -2 \end{pmatrix}, v_4 = \begin{pmatrix} -2 \\ -4 \\ 6 \end{pmatrix}$.
\end{enumerate}
\end{ejem}

\begin{thm}
Sea $A$ una matriz de $m \times n$. Entonces $\rho(A) + \nu (A) = n$.
\end{thm}




\begin{ejem}
\begin{enumerate}
    \item 
Sea $A = \begin{pmatrix} 1 & 2 & -1 \\ 2 & -1 & 3 \end{pmatrix}$. $\nu(A) = 1$, $\rho(A) = 2$. Luego $\rho(A) + \nu(A) = 3$
\\\

\item
Para $A = \begin{pmatrix} 1 & -1 & 3 \\ 2 & 0 & 4 \\ -1 & -3 & 1 \end{pmatrix}$. Calcule $\nu(A)$
\end{enumerate}
\end{ejem}

\begin{thm}
Sea $A$ una matriz de $n \times n$. Entonces $A$ es invertible si y solo si $\rho (A) = n$.
\end{thm}

\begin{thm}
El sistema $Ax = b$ tiene al menos una solución si y sólo si $b \in C_A$. Esto ocurre si y solo si $A$ y la matriz aumentada $(A,b)$ tienen el mismo rango.
\end{thm}





\begin{thm}
\textbf{(Resumen)} Sea $A$ una matriz de $n \times n$. Entonces las siguientes afirmaciones son equivalentes:

\begin{enumerate}
\item
$A$ es invertible.

\item
La única solución al sistema homogéneo $Ax = 0$ es la solución trivial.

\item
El sistema $Ax = b$ tiene una solución única para cada vector $b$ de dimensión $n$.

\item
$A$ es equivalente por filas a la matriz identidad $\mathbb{I}$ de $n \times n$.

\item
La forma escalonada por filas de $A$ tiene $n$ pivotes.

\item
Las columnas y filas de $A$ son linealmente independientes.

\item
$\det A \not=0$.

\item
$\nu(A) = 0$.

\item
$\rho(A) = n$.
\end{enumerate}
\end{thm}










